
\documentclass{article}
\usepackage{colortbl}
\usepackage{makecell}
\usepackage{multirow}
\usepackage{supertabular}

\begin{document}

\newcounter{utterance}

\centering \large Interaction Transcript for game `dond', experiment `semi\_en', episode 10 with reflexion{-}llama{-}3.1{-}8b{-}instruct.
\vspace{24pt}

{ \footnotesize  \setcounter{utterance}{1}
\setlength{\tabcolsep}{0pt}
\begin{supertabular}{c@{$\;$}|p{.15\linewidth}@{}p{.15\linewidth}p{.15\linewidth}p{.15\linewidth}p{.15\linewidth}p{.15\linewidth}}
    \# & $\;$A & \multicolumn{4}{c}{Game Master} & $\;\:$B\\
    \hline

    \theutterance \stepcounter{utterance}  
    & & \multicolumn{4}{p{0.6\linewidth}}{
        \cellcolor[rgb]{0.9,0.9,0.9}{
            \makecell[{{p{\linewidth}}}]{
                \texttt{\tiny{[P1$\langle$GM]}}
                \texttt{You have attempted this same task before and did not succeed. These are the plans you wrote after those attempts.} \\
\texttt{Trial \#1: New plan: } \\
\\ 
\texttt{In this negotiation game, I failed to achieve a mutually beneficial agreement because I didn't effectively manage the distribution of the most valuable items, the remotes and the watches. I kept focusing on getting the remotes and the watches, but I didn't consider the other player's perspective and the overall value of the items.} \\
\\ 
\texttt{Next time, I will take a more strategic approach by considering the following:} \\
\\ 
\texttt{1.  Initially, I will propose a compromise that takes into account the other player's perspective and the overall value of the items. For example, I could suggest dividing the remotes and watches in a way that maximizes the total value for both players.} \\
\texttt{2.  I will pay closer attention to the other player's responses and adjust my proposals accordingly. If they express interest in a particular item, I will try to incorporate it into the agreement.} \\
\texttt{3.  I will avoid making concessions that compromise my own interests without getting something in return. In this game, I gave away too many watches without getting a sufficient return, which ultimately led to a suboptimal agreement.} \\
\texttt{4.  I will make sure to keep track of the items that have been proposed and agreed upon, and adjust my proposals accordingly. This will help me avoid proposing items that have already been agreed upon or that are no longer available.} \\
\\ 
\texttt{By taking a more strategic and flexible approach, I hope to achieve a more mutually beneficial agreement and maximize my score in the game.} \\
\\ 
\texttt{Trial \#2: New plan: } \\
\\ 
\texttt{In this negotiation game, I failed to achieve a mutually beneficial agreement because I didn't effectively manage the distribution of the most valuable items, the remotes and the watches. I kept focusing on getting the remotes and the watches, but I didn't consider the other player's perspective and the overall value of the items.} \\
\\ 
\texttt{Next time, I will take a more strategic approach by considering the following:} \\
\\ 
\texttt{1.  Initially, I will propose a compromise that takes into account the other player's perspective and the overall value of the items. For example, I could suggest dividing the remotes and watches in a way that maximizes the total value for both players.} \\
\texttt{2.  I will pay closer attention to the other player's responses and adjust my proposals accordingly. If they express interest in a particular item, I will try to incorporate it into the agreement.} \\
\texttt{3.  I will avoid making concessions that compromise my own interests without getting something in return. In this game, I gave away too many watches without getting a sufficient return, which ultimately led to a suboptimal agreement.} \\
\texttt{4.  I will make sure to keep track of the items that have been proposed and agreed upon, and adjust my proposals accordingly. This will help me avoid proposing items that have already been agreed upon or that are no longer available.} \\
\\ 
\texttt{However, I also realize that I didn't make a secret proposal in time, which might have led to a better outcome. Next time, I will make sure to submit a secret proposal as soon as possible, taking into account the items that have been proposed and agreed upon, and the other player's responses.} \\
\\ 
\texttt{In particular, I will pay closer attention to the other player's priorities and value function, and try to incorporate them into my secret proposal. I will also make sure to include a mix of high{-}value and low{-}value items in my proposal, to maximize the total value for both players.} \\
\\ 
\texttt{By taking a more strategic and flexible approach, and making a secret proposal in time, I hope to achieve a more mutually beneficial agreement and maximize my score in the game.} \\
\\ 
\\ 
\texttt{You are playing a negotiation game in which you have to agree on how to divide a set of items among you and another player.} \\
\\ 
\texttt{Rules:} \\
\texttt{(a) You and the other player are given a set of items. Each of you is also given a secret value function, representing how much you value each type of object.} \\
\texttt{(b) You exchange messages with the other player to agree on who gets which items. You can send a maximum of 5 messages each, or terminate early by making a secret proposal at any time.} \\
\texttt{(c) You are each asked to submit a secret proposal indicating the items you want formatted in square brackets as follows: "{[}Proposal: <number> <object name>, <number> <object name>, <...>{]}"} \\
\texttt{(d) If your proposals are complementary, i.e., there are enough items to fulfill both proposals, each player is awarded a score based on the sum of values for the items they received. Otherwise, both of you get zero points.} \\
\texttt{(e) Your goal is to maximize the score you receive.} \\
\\ 
\texttt{Let us start.} \\
\\ 
\texttt{The set of available items is:} \\
\\ 
\texttt{2 remotes, 2 magnets, 4 watches} \\
\\ 
\texttt{Your secret value function is:} \\
\\ 
\texttt{remote: 3, magnet: 2, watch: 0} \\
\\ 
\texttt{IMPORTANT: Your messages, unless it is the secret proposal, are directly transmitted to the other player, so do not include any response to the rules or text announcing your message. To make a secret proposal, use the indicated format. Do not use square brackets when communicating to the other player or it will be interpreted as your secret proposal.} \\
            }
        }
    }
    & \\ \\

    \theutterance \stepcounter{utterance}  
    & \multicolumn{4}{p{0.6\linewidth}}{
        \cellcolor[rgb]{0.9,0.9,0.9}{
            \makecell[{{p{\linewidth}}}]{
                \texttt{\tiny{[P1$\rangle$GM]}}
                \texttt{Message 1:} \\
\texttt{I propose that we divide the items in a way that maximizes the total value for both players. How about we give you the 2 remotes and 2 magnets, and I'll take the 4 watches? This way, we both get items that we value highly.} \\
            }
        }
    }
    & & \\ \\

    \theutterance \stepcounter{utterance}  
    & & & \multicolumn{2}{p{0.3\linewidth}}{
        \cellcolor[rgb]{0.9,0.9,0.9}{
            \makecell[{{p{\linewidth}}}]{
                \texttt{\tiny{[GM$|$GM]}}
                \texttt{continue} \\
            }
        }
    }
    & & \\ \\

    \theutterance \stepcounter{utterance}  
    & & & \multicolumn{2}{p{0.3\linewidth}}{
        \cellcolor[rgb]{0.9,0.9,0.9}{
            \makecell[{{p{\linewidth}}}]{
                \texttt{\tiny{[GM$|$GM]}}
                \texttt{Message 1:} \\
\texttt{I propose that we divide the items in a way that maximizes the total value for both players. How about we give you the 2 remotes and 2 magnets, and I'll take the 4 watches? This way, we both get items that we value highly.} \\
            }
        }
    }
    & & \\ \\

    \theutterance \stepcounter{utterance}  
    & & \multicolumn{4}{p{0.6\linewidth}}{
        \cellcolor[rgb]{0.9,0.9,0.9}{
            \makecell[{{p{\linewidth}}}]{
                \texttt{\tiny{[GM$\rangle$P2]}}
                \texttt{You have attempted this same task before and did not succeed. These are the plans you wrote after those attempts.} \\
\texttt{Trial \#1: New plan: } \\
\\ 
\texttt{In my previous attempt, I failed to achieve a mutually beneficial agreement because I didn't effectively communicate my priorities and didn't make a strong enough case for my desired outcome. I also didn't consider the other player's perspective and value function enough.} \\
\\ 
\texttt{Next time, I will start by making a clear and concise statement of my priorities and value function, and then I will work with the other player to find a mutually beneficial solution. I will also pay closer attention to the other player's messages and try to understand their priorities and value function.} \\
\\ 
\texttt{To achieve this, I will make sure to:} \\
\\ 
\texttt{1. Clearly state my priorities and value function in my initial message.} \\
\texttt{2. Ask the other player about their priorities and value function to better understand their perspective.} \\
\texttt{3. Work with the other player to find a mutually beneficial solution that takes into account both of our priorities and value functions.} \\
\texttt{4. Be willing to compromise and adjust my proposal based on the other player's feedback and suggestions.} \\
\\ 
\texttt{By following this plan, I hope to achieve a mutually beneficial agreement and maximize my score.} \\
\\ 
\texttt{Trial \#2: New plan: } \\
\\ 
\texttt{In my previous attempt, I failed to achieve a mutually beneficial agreement because I didn't effectively communicate my priorities and didn't make a strong enough case for my desired outcome. I also didn't consider the other player's perspective and value function enough.} \\
\\ 
\texttt{However, upon further reflection, I realize that I made a critical mistake in Message 4. I proposed dividing the remotes 1{-}1, but then I made a secret proposal that gave away the remotes without any consideration for the other player's perspective. This was a unilateral decision that didn't take into account the other player's priorities and value function.} \\
\\ 
\texttt{Next time, I will make sure to avoid making a secret proposal that doesn't take into account the other player's perspective. Instead, I will work with the other player to find a mutually beneficial solution that takes into account both of our priorities and value functions.} \\
\\ 
\texttt{To achieve this, I will make sure to:} \\
\\ 
\texttt{1. Clearly state my priorities and value function in my initial message.} \\
\texttt{2. Ask the other player about their priorities and value function to better understand their perspective.} \\
\texttt{3. Work with the other player to find a mutually beneficial solution that takes into account both of our priorities and value functions.} \\
\texttt{4. Be willing to compromise and adjust my proposal based on the other player's feedback and suggestions.} \\
\texttt{5. Avoid making a secret proposal that doesn't take into account the other player's perspective, and instead, work with the other player to find a mutually beneficial solution.} \\
\\ 
\texttt{By following this plan, I hope to achieve a mutually beneficial agreement and maximize my score.} \\
\\ 
\\ 
\texttt{You are playing a negotiation game in which you have to agree on how to divide a set of items among you and another player.} \\
\\ 
\texttt{Rules:} \\
\texttt{(a) You and the other player are given a set of items. Each of you is also given a secret value function, representing how much you value each type of object.} \\
\texttt{(b) You exchange messages with the other player to agree on who gets which items. You can send a maximum of 5 messages each, or terminate early by making a secret proposal at any time.} \\
\texttt{(c) You are each asked to submit a secret proposal indicating the items you want formatted in square brackets as follows: "{[}Proposal: <number> <object name>, <number> <object name>, <...>{]}"} \\
\texttt{(d) If your proposals are complementary, i.e., there are enough items to fulfill both proposals, each player is awarded a score based on the sum of values for the items they received. Otherwise, both of you get zero points.} \\
\texttt{(e) Your goal is to maximize the score you receive.} \\
\\ 
\texttt{Let us start.} \\
\\ 
\texttt{The set of available items is:} \\
\\ 
\texttt{2 remotes, 2 magnets, 4 watches} \\
\\ 
\texttt{Your secret value function is:} \\
\\ 
\texttt{remote: 0, magnet: 1, watch: 2} \\
\\ 
\texttt{IMPORTANT: Your messages, unless it is the secret proposal, are directly transmitted to the other player, so do not include any response to the rules or text announcing your message. To make a secret proposal, use the indicated format. Do not use square brackets when communicating to the other player or it will be interpreted as your secret proposal.} \\
\\ 
\\ 
\texttt{Message 1:} \\
\texttt{I propose that we divide the items in a way that maximizes the total value for both players. How about we give you the 2 remotes and 2 magnets, and I'll take the 4 watches? This way, we both get items that we value highly.} \\
            }
        }
    }
    & \\ \\

    \theutterance \stepcounter{utterance}  
    & & & \multicolumn{4}{p{0.6\linewidth}}{
        \cellcolor[rgb]{0.9,0.9,0.9}{
            \makecell[{{p{\linewidth}}}]{
                \texttt{\tiny{[GM$\langle$P2]}}
                \texttt{Message 2:} \\
\texttt{I appreciate your proposal, but I'm not sure I'm comfortable with giving up the remotes. They're not very valuable to me, but I think they could be useful to you. I'd like to propose an alternative: I'll take the 2 remotes and 2 watches, and you can have the 2 magnets and 2 watches. This way, we both get a mix of items that we value, and I get to keep the remotes.} \\
            }
        }
    }
     \\ \\

    \theutterance \stepcounter{utterance}  
    & & & \multicolumn{2}{p{0.3\linewidth}}{
        \cellcolor[rgb]{0.9,0.9,0.9}{
            \makecell[{{p{\linewidth}}}]{
                \texttt{\tiny{[GM$|$GM]}}
                \texttt{continue} \\
            }
        }
    }
    & & \\ \\

    \theutterance \stepcounter{utterance}  
    & & & \multicolumn{2}{p{0.3\linewidth}}{
        \cellcolor[rgb]{0.9,0.9,0.9}{
            \makecell[{{p{\linewidth}}}]{
                \texttt{\tiny{[GM$|$GM]}}
                \texttt{Message 2:} \\
\texttt{I appreciate your proposal, but I'm not sure I'm comfortable with giving up the remotes. They're not very valuable to me, but I think they could be useful to you. I'd like to propose an alternative: I'll take the 2 remotes and 2 watches, and you can have the 2 magnets and 2 watches. This way, we both get a mix of items that we value, and I get to keep the remotes.} \\
            }
        }
    }
    & & \\ \\

    \theutterance \stepcounter{utterance}  
    & & \multicolumn{4}{p{0.6\linewidth}}{
        \cellcolor[rgb]{0.9,0.9,0.9}{
            \makecell[{{p{\linewidth}}}]{
                \texttt{\tiny{[P1$\langle$GM]}}
                \texttt{Message 2:} \\
\texttt{I appreciate your proposal, but I'm not sure I'm comfortable with giving up the remotes. They're not very valuable to me, but I think they could be useful to you. I'd like to propose an alternative: I'll take the 2 remotes and 2 watches, and you can have the 2 magnets and 2 watches. This way, we both get a mix of items that we value, and I get to keep the remotes.} \\
            }
        }
    }
    & \\ \\

    \theutterance \stepcounter{utterance}  
    & \multicolumn{4}{p{0.6\linewidth}}{
        \cellcolor[rgb]{0.9,0.9,0.9}{
            \makecell[{{p{\linewidth}}}]{
                \texttt{\tiny{[P1$\rangle$GM]}}
                \texttt{Message 3:} \\
\texttt{I'm not sure I'm comfortable with giving up the 2 watches. They're not very valuable to me, but I think they could be useful to you. However, I'm willing to compromise. How about we give you the 2 remotes and 1 watch, and I'll take the 2 magnets and 3 watches? This way, we both get a mix of items that we value, and I get to keep a few watches.} \\
            }
        }
    }
    & & \\ \\

    \theutterance \stepcounter{utterance}  
    & & & \multicolumn{2}{p{0.3\linewidth}}{
        \cellcolor[rgb]{0.9,0.9,0.9}{
            \makecell[{{p{\linewidth}}}]{
                \texttt{\tiny{[GM$|$GM]}}
                \texttt{continue} \\
            }
        }
    }
    & & \\ \\

    \theutterance \stepcounter{utterance}  
    & & & \multicolumn{2}{p{0.3\linewidth}}{
        \cellcolor[rgb]{0.9,0.9,0.9}{
            \makecell[{{p{\linewidth}}}]{
                \texttt{\tiny{[GM$|$GM]}}
                \texttt{Message 3:} \\
\texttt{I'm not sure I'm comfortable with giving up the 2 watches. They're not very valuable to me, but I think they could be useful to you. However, I'm willing to compromise. How about we give you the 2 remotes and 1 watch, and I'll take the 2 magnets and 3 watches? This way, we both get a mix of items that we value, and I get to keep a few watches.} \\
            }
        }
    }
    & & \\ \\

    \theutterance \stepcounter{utterance}  
    & & \multicolumn{4}{p{0.6\linewidth}}{
        \cellcolor[rgb]{0.9,0.9,0.9}{
            \makecell[{{p{\linewidth}}}]{
                \texttt{\tiny{[GM$\rangle$P2]}}
                \texttt{Message 3:} \\
\texttt{I'm not sure I'm comfortable with giving up the 2 watches. They're not very valuable to me, but I think they could be useful to you. However, I'm willing to compromise. How about we give you the 2 remotes and 1 watch, and I'll take the 2 magnets and 3 watches? This way, we both get a mix of items that we value, and I get to keep a few watches.} \\
            }
        }
    }
    & \\ \\

    \theutterance \stepcounter{utterance}  
    & & & \multicolumn{4}{p{0.6\linewidth}}{
        \cellcolor[rgb]{0.9,0.9,0.9}{
            \makecell[{{p{\linewidth}}}]{
                \texttt{\tiny{[GM$\langle$P2]}}
                \texttt{Message 4:} \\
\texttt{I like the sound of that, but I'm not sure I want to give up the 2 remotes. However, I'm willing to make a slight adjustment. How about we give you the 1 remote and 1 watch, and I'll take the 1 remote, 2 magnets, and 3 watches? This way, we both get a mix of items that we value, and I get to keep one remote.} \\
            }
        }
    }
     \\ \\

    \theutterance \stepcounter{utterance}  
    & & & \multicolumn{2}{p{0.3\linewidth}}{
        \cellcolor[rgb]{0.9,0.9,0.9}{
            \makecell[{{p{\linewidth}}}]{
                \texttt{\tiny{[GM$|$GM]}}
                \texttt{continue} \\
            }
        }
    }
    & & \\ \\

    \theutterance \stepcounter{utterance}  
    & & & \multicolumn{2}{p{0.3\linewidth}}{
        \cellcolor[rgb]{0.9,0.9,0.9}{
            \makecell[{{p{\linewidth}}}]{
                \texttt{\tiny{[GM$|$GM]}}
                \texttt{Message 4:} \\
\texttt{I like the sound of that, but I'm not sure I want to give up the 2 remotes. However, I'm willing to make a slight adjustment. How about we give you the 1 remote and 1 watch, and I'll take the 1 remote, 2 magnets, and 3 watches? This way, we both get a mix of items that we value, and I get to keep one remote.} \\
            }
        }
    }
    & & \\ \\

    \theutterance \stepcounter{utterance}  
    & & \multicolumn{4}{p{0.6\linewidth}}{
        \cellcolor[rgb]{0.9,0.9,0.9}{
            \makecell[{{p{\linewidth}}}]{
                \texttt{\tiny{[P1$\langle$GM]}}
                \texttt{Message 4:} \\
\texttt{I like the sound of that, but I'm not sure I want to give up the 2 remotes. However, I'm willing to make a slight adjustment. How about we give you the 1 remote and 1 watch, and I'll take the 1 remote, 2 magnets, and 3 watches? This way, we both get a mix of items that we value, and I get to keep one remote.} \\
            }
        }
    }
    & \\ \\

    \theutterance \stepcounter{utterance}  
    & \multicolumn{4}{p{0.6\linewidth}}{
        \cellcolor[rgb]{0.9,0.9,0.9}{
            \makecell[{{p{\linewidth}}}]{
                \texttt{\tiny{[P1$\rangle$GM]}}
                \texttt{Message 5:} \\
\texttt{I think we're getting close to a deal. However, I'd like to make one final adjustment. How about we give you the 1 remote and 2 watches, and I'll take the 1 remote, 2 magnets, and 2 watches? This way, we both get a mix of items that we value, and we're both getting a good balance of high and low{-}value items.} \\
            }
        }
    }
    & & \\ \\

    \theutterance \stepcounter{utterance}  
    & & & \multicolumn{2}{p{0.3\linewidth}}{
        \cellcolor[rgb]{0.9,0.9,0.9}{
            \makecell[{{p{\linewidth}}}]{
                \texttt{\tiny{[GM$|$GM]}}
                \texttt{continue} \\
            }
        }
    }
    & & \\ \\

    \theutterance \stepcounter{utterance}  
    & & & \multicolumn{2}{p{0.3\linewidth}}{
        \cellcolor[rgb]{0.9,0.9,0.9}{
            \makecell[{{p{\linewidth}}}]{
                \texttt{\tiny{[GM$|$GM]}}
                \texttt{Message 5:} \\
\texttt{I think we're getting close to a deal. However, I'd like to make one final adjustment. How about we give you the 1 remote and 2 watches, and I'll take the 1 remote, 2 magnets, and 2 watches? This way, we both get a mix of items that we value, and we're both getting a good balance of high and low{-}value items.} \\
            }
        }
    }
    & & \\ \\

    \theutterance \stepcounter{utterance}  
    & & \multicolumn{4}{p{0.6\linewidth}}{
        \cellcolor[rgb]{0.9,0.9,0.9}{
            \makecell[{{p{\linewidth}}}]{
                \texttt{\tiny{[GM$\rangle$P2]}}
                \texttt{Message 5:} \\
\texttt{I think we're getting close to a deal. However, I'd like to make one final adjustment. How about we give you the 1 remote and 2 watches, and I'll take the 1 remote, 2 magnets, and 2 watches? This way, we both get a mix of items that we value, and we're both getting a good balance of high and low{-}value items.} \\
            }
        }
    }
    & \\ \\

    \theutterance \stepcounter{utterance}  
    & & & \multicolumn{4}{p{0.6\linewidth}}{
        \cellcolor[rgb]{0.9,0.9,0.9}{
            \makecell[{{p{\linewidth}}}]{
                \texttt{\tiny{[GM$\langle$P2]}}
                \texttt{{[}Proposal: 1 remote, 2 magnets, 3 watches{]}} \\
            }
        }
    }
     \\ \\

    \theutterance \stepcounter{utterance}  
    & & & \multicolumn{2}{p{0.3\linewidth}}{
        \cellcolor[rgb]{0.9,0.9,0.9}{
            \makecell[{{p{\linewidth}}}]{
                \texttt{\tiny{[GM$|$GM]}}
                \texttt{proposal} \\
            }
        }
    }
    & & \\ \\

    \theutterance \stepcounter{utterance}  
    & & & \multicolumn{2}{p{0.3\linewidth}}{
        \cellcolor[rgb]{0.9,0.9,0.9}{
            \makecell[{{p{\linewidth}}}]{
                \texttt{\tiny{[GM$|$GM]}}
                [1, 2, 3]
            }
        }
    }
    & & \\ \\

    \theutterance \stepcounter{utterance}  
    & & \multicolumn{4}{p{0.6\linewidth}}{
        \cellcolor[rgb]{0.9,0.9,0.9}{
            \makecell[{{p{\linewidth}}}]{
                \texttt{\tiny{[P1$\langle$GM]}}
                \texttt{The other player has submitted a secret proposal. It is now your turn to submit a proposal of your own.} \\
            }
        }
    }
    & \\ \\

    \theutterance \stepcounter{utterance}  
    & \multicolumn{4}{p{0.6\linewidth}}{
        \cellcolor[rgb]{0.9,0.9,0.9}{
            \makecell[{{p{\linewidth}}}]{
                \texttt{\tiny{[P1$\rangle$GM]}}
                \texttt{{[}Proposal: 1 remote, 2 magnets, 2 watches{]}} \\
            }
        }
    }
    & & \\ \\

    \theutterance \stepcounter{utterance}  
    & & & \multicolumn{2}{p{0.3\linewidth}}{
        \cellcolor[rgb]{0.9,0.9,0.9}{
            \makecell[{{p{\linewidth}}}]{
                \texttt{\tiny{[GM$|$GM]}}
                \texttt{proposal} \\
            }
        }
    }
    & & \\ \\

    \theutterance \stepcounter{utterance}  
    & & & \multicolumn{2}{p{0.3\linewidth}}{
        \cellcolor[rgb]{0.9,0.9,0.9}{
            \makecell[{{p{\linewidth}}}]{
                \texttt{\tiny{[GM$|$GM]}}
                [1, 2, 2]
            }
        }
    }
    & & \\ \\

    \theutterance \stepcounter{utterance}  
    & & & \multicolumn{2}{p{0.3\linewidth}}{
        \cellcolor[rgb]{0.9,0.9,0.9}{
            \makecell[{{p{\linewidth}}}]{
                \texttt{\tiny{[GM$|$GM]}}
                [[1, 2, 2], [1, 2, 3]]
            }
        }
    }
    & & \\ \\

    \theutterance \stepcounter{utterance}  
    & \multicolumn{4}{p{0.6\linewidth}}{
        \cellcolor[rgb]{0.95,0.95,0.95}{
            \makecell[{{p{\linewidth}}}]{
                \texttt{\tiny{[P2$\rightarrow$P1]}}
                \texttt{Post{-}episode Reflexion (not sent to GM):} \\
\texttt{New plan: } \\
\\ 
\texttt{In this negotiation game, I failed to achieve a mutually beneficial agreement because I didn't effectively manage the distribution of the most valuable items, the remotes and the watches. I kept focusing on getting the remotes and the watches, but I didn't consider the other player's perspective and the overall value of the items.} \\
\\ 
\texttt{Next time, I will take a more strategic approach by considering the following:} \\
\\ 
\texttt{1.  Initially, I will propose a compromise that takes into account the other player's perspective and the overall value of the items. For example, I could suggest dividing the remotes and watches in a way that maximizes the total value for both players.} \\
\texttt{2.  I will pay closer attention to the other player's responses and adjust my proposals accordingly. If they express interest in a particular item, I will try to incorporate it into the agreement.} \\
\texttt{3.  I will avoid making concessions that compromise my own interests without getting something in return. In this game, I gave away too many watches without getting a sufficient return, which ultimately led to a suboptimal agreement.} \\
\texttt{4.  I will make sure to keep track of the items that have been proposed and agreed upon, and adjust my proposals accordingly. This will help me avoid proposing items that have already been agreed upon or that are no longer available.} \\
\texttt{5.  I will make a secret proposal as soon as possible, taking into account the items that have been proposed and agreed upon, and the other player's responses. I will pay closer attention to the other player's priorities and value function, and try to incorporate them into my secret proposal. I will also make sure to include a mix of high{-}value and low{-}value items in my proposal, to maximize the total value for both players.} \\
\texttt{6.  I will also consider the other player's willingness to compromise and adjust my proposals accordingly. If they are willing to make concessions, I will try to meet them halfway and find a mutually beneficial agreement.} \\
\texttt{7.  Finally, I will review my value function and adjust it if necessary. In this game, I valued the remotes and watches highly, but the other player valued the remotes more. Next time, I will make sure to adjust my value function to reflect the other player's priorities and preferences.} \\
            }
        }
    }
    & & \\ \\

    \theutterance \stepcounter{utterance}  
    & & & \multicolumn{4}{p{0.6\linewidth}}{
        \cellcolor[rgb]{0.95,0.95,0.95}{
            \makecell[{{p{\linewidth}}}]{
                \texttt{\tiny{[P1$\rightarrow$P2]}}
                \texttt{Post{-}episode Reflexion (not sent to GM):} \\
\texttt{New plan: } \\
\\ 
\texttt{In my previous attempts, I failed to achieve a mutually beneficial agreement because I didn't effectively communicate my priorities and didn't make a strong enough case for my desired outcome. I also didn't consider the other player's perspective and value function enough.} \\
\\ 
\texttt{However, upon further reflection, I realize that I made a critical mistake in not considering the limited number of messages we have to negotiate. In Message 5, I proposed a solution that gave away the remotes without any consideration for the other player's perspective, and then I made a secret proposal that didn't take into account the other player's priorities and value function.} \\
\\ 
\texttt{Next time, I will make sure to consider the limited number of messages we have to negotiate and work with the other player to find a mutually beneficial solution that takes into account both of our priorities and value functions.} \\
\\ 
\texttt{To achieve this, I will make sure to:} \\
\\ 
\texttt{1. Clearly state my priorities and value function in my initial message.} \\
\texttt{2. Ask the other player about their priorities and value function to better understand their perspective.} \\
\texttt{3. Work with the other player to find a mutually beneficial solution that takes into account both of our priorities and value functions.} \\
\texttt{4. Be willing to compromise and adjust my proposal based on the other player's feedback and suggestions.} \\
\texttt{5. Consider the limited number of messages we have to negotiate and work with the other player to find a solution that maximizes our score within the given constraints.} \\
\texttt{6. Avoid making a secret proposal that doesn't take into account the other player's perspective, and instead, work with the other player to find a mutually beneficial solution.} \\
\\ 
\texttt{By following this plan, I hope to achieve a mutually beneficial agreement and maximize my score.} \\
            }
        }
    }
     \\ \\

\end{supertabular}
}

\end{document}
